\documentclass[10pt,a4paper]{amsart}
\usepackage[margin=19mm]{geometry}
\usepackage{amsmath,amssymb,mathtools}
\usepackage[scr=rsfs,cal=euler]{mathalfa}
\usepackage{xfrac}
\usepackage[unicode,hidelinks,bookmarks=false]{hyperref}
\newcommand{\ie}{i.e.\ }
\newcommand{\cf}{cf.\ }
\newcommand{\resp}{respectively\ }
\newcommand{\Subsection}{\subsection}
\providecommand{\nobreakdash}{\nobreak\hbox{-}}
\setlength{\emergencystretch}{2em}
\multlinegap0pt
\usepackage{amssymb}
\usepackage{float}
\newtheorem{proposition}{Proposition}
\newcommand{\Reff}{R}
\title{Computation of Unknotting Numbers:\\ \mbox{Which Knot Breaks the Bernhard--Jablan Conjecture?}}
\author{Seong-Jin Lee}
\date{Source: arXiv:2609.09861v2 (CC BY 4.0)}
\begin{document}
\maketitle

\noindent\emph{Source and licence.} Computation of Unknotting Numbers:\\ \mbox{Which Knot Breaks the Bernhard--Jablan Conjecture?}. Version arXiv:2609.09861v2 (CC BY 4.0), distributed by arXiv under the Creative Commons Attribution 4.0 International licence, http://creativecommons.org/licenses/by/4.0/, as recorded in the arXiv licence metadata for this exact version and shown on the abstract page https://arxiv.org/abs/2609.09861v2. Full licence notice: \texttt{licenses/arxiv-2609.09861-CC-BY-4.0.txt}.\par\smallskip

\noindent\textit{Editorial scope.} Counted unit: the article's proposition prop:sigma (source0.tex lines 1336--1345). The article's complete original proof of it is delivered with it (source0.tex lines 1347--1360, one \texttt{proof} environment of 14 lines). No mathematics is written, repaired or re-derived: the statement and the proof are the article's own lines, in the article's own notation, and no retained line is substituted or re-flowed.  Retained uncounted as the source-local context the proof needs: lines 1330--1333.  Nothing was omitted from any retained span, so the package carries no omission record.  No retained line contains a citation, so the wrapper carries no bibliography and no bibliography entry.  No retained line contains a figure environment, an image-inclusion command or drawing code, so no image is delivered and the package declares no asset dependency.  Editorial formatting only, in addition: an AMS \texttt{amsart} preamble that restates the article's own declarations and macros used by the retained lines, namely the environments \texttt{proposition} on separate counters, together with the standard packages those lines need.  The retained environments print in this document as Proposition 1 in order of appearance, whereas the article prints the same items with the numbering of its own full document.  The complete native source of the article as distributed in the arXiv e-print for this exact version is bundled unchanged as \texttt{sources/209856/source0.tex}.
\begin{equation}\label{eq:gl}
    \sigma(K)=\operatorname{sign}(G)-\mu(F),\qquad
    \mu(F)=\sum_{c\text{ of type II}}\eta(c).
\end{equation}

\begin{proposition}\label{prop:sigma}
Let $c$ be a crossing of $D$, let $e$ be its Tait edge, and let $K_c$ be the knot after
the crossing change. Then
\begin{equation}\label{eq:det-change}
    \det K_c=\det K\,|1-2\Reff(e)|.
\end{equation}
For a type I crossing, the signature decreases by two if $\Reff(e)>1/2$ and is unchanged
if $\Reff(e)<1/2$. For a type II crossing, it increases by two if $\Reff(e)<1/2$ and
is unchanged if $\Reff(e)>1/2$. The value $\Reff(e)=1/2$ cannot occur.
\end{proposition}

\begin{proof}
If $x$ is the difference of the coordinate vectors of the two regions incident to $e$,
with the deleted coordinate set to zero, the changed Goeritz matrix is
$G'=G-2xx^T$. The matrix determinant lemma gives
\[
    \det G'=\det G(1-2x^TG^{-1}x),\qquad x^TG^{-1}x=\Reff(e).
\]
This proves \eqref{eq:det-change}. The determinant of a knot is odd, so $G'$ is
nonsingular and $\Reff(e)\ne1/2$. As $G$ is positive definite, the rank-one change either
preserves its signature or decreases it by two, according to the sign of
$1-2\Reff(e)$. The type of the crossing is unchanged, while its Goeritz sign reverses.
Consequently $\mu$ decreases by two at a type II crossing and is unchanged at a type I
crossing. Substitution in \eqref{eq:gl} gives the signature statements.
\end{proof}

\end{document}
